\section{Overlap conventions and dimension of the attractor}
\label{sec:corollaries}

We first reconcile the two usual word-length conventions for an exact overlap.
The following observation also appears in \cite[footnote~3]{Hochman2014}.
The nonzero contraction hypothesis is useful here as well as in the proof of
the measure theorem.

\begin{lemma}\label{lem:overlap-lengths}
For a finite family of maps $x\mapsto r_i x+t_i$ with $0<|r_i|<1$,
an equality of maps for two distinct nonempty words of arbitrary lengths
implies such an equality for two distinct words of the same length.
\end{lemma}
\begin{proof}
Suppose $\varphi_u=\varphi_v$ and $u\ne v$. If $u$ were a proper prefix
of $v$, writing $v=uw$ and cancelling the invertible map $\varphi_u$
would give $\varphi_w=\mathrm{id}$. This is impossible because a nonempty
word has absolute contraction less than one. The same applies with the
words interchanged. Now
$\varphi_{uv}=\varphi_u\circ\varphi_v
=\varphi_v\circ\varphi_u=\varphi_{vu}$.
The words $uv$ and $vu$ have the same length and are distinct: if they
were equal, their first $\min\{|u|,|v|\}$ symbols would show that either
$u=v$ or the shorter word is a proper prefix of the longer.
\end{proof}

\begin{corollary}[No exact overlaps]\label{cor:no-overlaps}
Let $\Phi$ be a finite nonempty indexed family of real similarities
with $0<|r_i|<1$ and no exact overlaps.  For every strictly positive
probability vector $p$,
\[
 \dim_H\mu_{\Phi,p}=\min\{1,H(p)/\chi(\Phi,p)\}.
\]
\end{corollary}
\begin{proof}
At every length $n$, distinct words give distinct complete maps.
Hence $G_n$ is an injective function of its symbol word, so
$H(G_n)=nH(p)$ and $h_{\mathrm{RW}}=H(p)$.
Apply Theorem~\ref{thm:main}.
\end{proof}

Let $K_\Phi$ be the compact attractor, equivalently the set of all coding
limits. Its similarity dimension $s_*$ is the unique number $s_*\ge0$
satisfying
\[
 \sum_{i\in\Lambda}|r_i|^{s_*}=1.
\]
For at least two symbols, existence and uniqueness follow because this sum
is continuous and strictly decreasing from $|\Lambda|>1$ to zero.
For one symbol its unique solution is zero.

\begin{corollary}\label{cor:set}
For a finite nonempty family of similarities of $\R$ with arbitrary real
translations and $0<|r_i|<1$, absence of exact overlaps implies
\[
 \dim_H K_\Phi=\min\{1,s_*\}.
\]
\end{corollary}
\begin{proof}
For one map the attractor is its fixed point and $s_*=0$.
For at least two maps choose the strictly positive weights
$p_i=|r_i|^{s_*}$. Then $H(p)=s_*\chi(\Phi,p)$.
Corollary~\ref{cor:no-overlaps} gives a measure supported on $K_\Phi$ of
dimension $\min\{1,s_*\}$, which is a lower bound for its set dimension.
For the upper bound, if $a>s_*$, the level-$n$ cylinder sets cover
$K_\Phi$, their diameters tend uniformly to zero, and
\[
 \sum_{u\in\Lambda^n}(\operatorname{diam}\varphi_u(K_\Phi))^a
 = (\operatorname{diam}K_\Phi)^a
   \left(\sum_i|r_i|^a\right)^n\longrightarrow0.
\]
The zero-diameter case is immediate. Thus $\dim_H K_\Phi\le s_*$,
and the ambient upper bound is one.
\end{proof}

\begin{corollary}\label{cor:homogeneous}
Let $m\ge2$, $0<\lambda<1$, and $t_1,\ldots,t_m\in\R$.
If the maps $x\mapsto\lambda x+t_i$ have no exact overlaps, then for
every strictly positive probability vector $p$ their self-similar measure has
\[
 \dim_H\mu=\min\left\{1,\frac{H(p)}{\log(1/\lambda)}\right\}.
\]
\end{corollary}
\begin{proof}
Apply Corollary~\ref{cor:no-overlaps} with $r_i=\lambda$ for every $i$.
Since the weights sum to one, $\chi=\log(1/\lambda)$.
\end{proof}

\begin{corollary}[Three maps]\label{cor:three-maps}
Let $0<\lambda<1$ and $t\in\R$.  If the indexed maps
\[
 x\longmapsto\lambda x,\qquad
 x\longmapsto\lambda x+1,\qquad
 x\longmapsto\lambda x+t
\]
have no exact overlaps, then for every $p_0,p_1,p_2>0$ with
$p_0+p_1+p_2=1$, their self-similar measure satisfies
\[
 \dim_H\mu
 =\min\left\{1,\frac{-\sum_{i=0}^2 p_i\log p_i}
                         {\log(1/\lambda)}\right\}.
\]
\end{corollary}
\begin{proof}
Apply Corollary~\ref{cor:homogeneous} with translations $0,1,t$.
\end{proof}

No claim of absolute continuity, effective separation, or a quantitative rate
of convergence follows from these dimension statements. In particular, the
finite upper endpoints of the scale ranges used in the proof have no asserted
effective bound.
